\section{夏季星空與天體}
夏季大三角：天鷹座的河鼓二（天狼星）、天鵝座的天津四、天琴座的織女星

\subsection{天蠍座 Scorpius}
\begin{enumerate}
    \item 黃道帶星座之一
    \item 最亮星：心宿二（天蠍座\(\alpha\)）(紅色)，視星等為0.96，是全天第十五亮星
\end{enumerate}
天蠍座中的深空天體：
\newline

% {
%     \centering
%     \begin{tabular}{c | c | c | c | c | c}
%         編號 & 種類 & 別名 & 視尺寸 & 拍攝技巧 & 曝光時間\\
%         \hline
%         NGC 6334 & 發射星雲 & 貓掌星雲 & \(\ang{;40;}\times\ang{;30;}\) & RGB/SHO & > 10 hr \\
%         IC 4628 & 發射星雲 & 斑節蝦星雲 & \(\ang{1;30;}\) & RGB/SHO & > 6 hr \\
%         M4 & 球狀星團 & - & \(\ang{;36;}\) & RGB & - \\
%         M6/NGC 6405 & 球狀星團 & 蝴蝶星團 & \(\ang{;10;}\) & RGB & - \\
%         M7/NGC 6475 & 疏散星團 & 托勒密星團 & \(\ang{1;20;}\) & RGB & - \\
%         M80/NGC 6093 & 球狀星團 & - & \(\ang{;10;}\) & RGB & - \\
%         Sh2-1 + Sh2-7 & 發射/反射星雲 & - & \(\ang{6;;}\times\ang{4;;}\) & RGB & - \\
%         Sh2-3/RCW 120 & 發射星雲 & 綠環星雲 & \(\ang{;6;}\) & RGB/SHO & > 8 hr \\
%         Sh2-11/NGC 6357 & 發射星雲 & 戰爭與和平星雲 & \(\ang{50;;}\times\ang{;40;}\) & RGB/SHO & > 6 hr \\
%     \end{tabular}\par
% }

\begin{center}
\begin{longtable}{c | c | c | c | c | c}
    % \hline
    編號 & 種類 & 別名 & 視尺寸 & RA & Dec \\
    \hline
    \endfirsthead

    % \hline
    % 編號 & 種類 & 別名 & 視尺寸 & 拍攝技巧 & 曝光時間 \\
    % \hline
    % \endhead

    % \hline
    % \multicolumn{6}{r}{接下頁} \\
    \endfoot

    % \hline
    \endlastfoot

    NGC 6334 & 發射星雲 & 貓掌星雲 & \(\ang{;40;} \times \ang{;30;}\) & \(17^\text{h }20^\text{m }49^\text{s}\) & \(-\ang{36;6;10}\) \\ % RGB/SHO & > 10 hr \\
    IC 4628 & 發射星雲 & 斑節蝦星雲 & \(\ang{1;30;}\) & \(16^\text{h }56^\text{m }57^\text{s}\) & \(-\ang{40;27;3}\) \\ % RGB/SHO & > 6 hr \\
    M4 & 球狀星團 & The Cat's Eye & \(\ang{;36;}\) & \(16^\text{h }23^\text{m }35^\text{s}\) & \(-\ang{26;31;32}\) \\ % RGB & - \\
    M6/NGC 6405 & 球狀星團 & 蝴蝶星團 & \(\ang{;10;}\) & \(17^\text{h }40^\text{m }20^\text{s}\) & \(-\ang{32;15;15}\) \\ % RGB & - \\
    M7/NGC 6475 & 疏散星團 & 托勒密星團 & \(\ang{1;20;}\) & \(17^\text{h }53^\text{m }51^\text{s}\) & \(-\ang{34;47;34}\) \\ % RGB & - \\
    M80/NGC 6093 & 球狀星團 & - & \(\ang{;10;}\) & \(16^\text{h }17^\text{m }03^\text{s}\) & \(-\ang{22;58;30}\) \\ % RGB & - \\
    Sh2-1 + Sh2-7 & 發射/反射星雲 & - & \(\ang{6;;} \times \ang{4;;}\) & \(15^\text{h }58^\text{m }51^\text{s}\) & \(-\ang{26;7;14}\) \\ % RGB & - \\
    Sh2-3/RCW 120 & 發射星雲 & 綠環星雲 & \(\ang{;6;}\) & \(17^\text{h }12^\text{m }23^\text{s}\) & \(-\ang{38;26;51}\) \\ % RGB/SHO & > 8 hr \\
    Sh2-11/NGC 6357 & 發射星雲 & 戰爭與和平星雲 & \(\ang{50;;} \times \ang{;40;}\) & \(17^\text{h }24^\text{m }43^\text{s}\) & \(-\ang{34;12;5}\) \\ % RGB/SHO & > 6 hr \\

\end{longtable}
\end{center}


\subsection{人馬座 Sagittarius}
\begin{enumerate}
    \item 最亮星：箕宿三，視星等1.8
    \item 可由天鵝座的天津四向天鷹座的牛郎星延長一倍找到
\end{enumerate}
人馬座中的深空天體：
\newline

\begin{center}
\begin{longtable}{c | c | c | c | c | c}
    % \hline
    編號 & 種類 & 別名 & 視尺寸 & RA & Dec \\
    \hline
    \endfirsthead

    % \hline
    % 編號 & 種類 & 別名 & 視尺寸 & 拍攝技巧 & 曝光時間 \\
    % \hline
    % \endhead

    % \hline
    % \multicolumn{6}{r}{接下頁} \\
    \endfoot

    % \hline
    \endlastfoot

    M8/NGC 6523 & 發射星雲 & 礁湖星雲 & \(\ang{1;30;} \times \ang{;40;}\) & \(18^\text{h }03^\text{m }41^\text{s}\) & \(-\ang{24;22;49}\) \\ % RGB/SHO & > 4 hr \\
    M17/NGC 6618 & 發射星雲 & 歐米茄星雲 & \(\ang{;46;} \times \ang{;37;}\) & \(18^\text{h }20^\text{m }47^\text{s}\) & \(-\ang{16;10;18}\) \\ % RGB/SHO & > 4 hr \\
    M18/NGC 6613 & 疏散星團 & - & \(\ang{;9;48}\) & \(18^\text{h }19^\text{m }58^\text{s}\) & \(-\ang{17;6;6}\) \\ % RGB & - \\
    M20/NGC 6514 & 發射/反射/暗星雲 & 三裂星雲 & \(\ang{;28;}\) & \(18^\text{h }02^\text{m }42^\text{s}\) & \(-\ang{22;58;18}\) \\ % RGB & > 4 hr \\
    M21/NGC 6531 & 疏散星團 & - & \(\ang{;14;}\) & \(18^\text{h }04^\text{m }13^\text{s}\) & \(-\ang{22;29;23}\) \\ % RGB & - \\
    M22/NGC 6656 & 球狀星團 & - & \(\ang{;32;}\) & \(18^\text{h }36^\text{m }24^\text{s}\) & \(-\ang{23;54;11}\) \\ % RGB & - \\
    M23/NGC 6494 & 疏散星團 & - & \(\ang{;35;}\) & \(17^\text{h }57^\text{m }05^\text{s}\) & \(-\ang{18;59;7}\) \\ % RGB & - \\
    M24/IC 4715 & 恆星雲 & 小人馬座恆星雲 & \(\ang{2;;} \times \ang{1;;}\) & \(18^\text{h }16^\text{m }55^\text{s}\) & \(-\ang{18;30;51}\) \\ % RGB & - \\
    M25/IC 4725 & 疏散星團 & - & \(\ang{;36;}\) & \(18^\text{h }31^\text{m }46^\text{s}\) & \(-\ang{19;6;53}\) \\ % RGB & - \\
    M28/NGC 6626 & 球狀星團 & - & \(\ang{;14;}\) & \(18^\text{h }24^\text{m }33^\text{s}\) & \(-\ang{24;52;10}\) \\ % RGB & - \\
    M54/NGC 6715 & 球狀星團 & - & \(\ang{;12;}\) & \(18^\text{h }55^\text{m }03^\text{s}\) & \(-\ang{30;28;42}\) \\ % RGB & - \\
    M55/NGC 6809 & 球狀星團 & - & \(\ang{;19;}\) & \(19^\text{h }39^\text{m }59^\text{s}\) & \(-\ang{30;57;43}\) \\ % RGB & - \\
    M69/NGC 6637 & 球狀星團 & - & \(\ang{;10;48}\) & \(18^\text{h }31^\text{m }23^\text{s}\) & \(-\ang{32;20;53}\) \\ % RGB & - \\
    M70/NGC 6681 & 球狀星團 & - & \(\ang{;8;}\) & \(18^\text{h }43^\text{m }13^\text{s}\) & \(-\ang{32;17;30}\) \\ % RGB & - \\
    M75/NGC 6864 & 球狀星團 & - & \(\ang{;6;48}\) & \(20^\text{h }06^\text{m }05^\text{s}\) & \(-\ang{21;55;19}\) \\ % RGB & - \\
    \makecell{IC 1284/Sh2-37 \\ NGC 6589 \\ NGC 6590 \\ LDN 315} & \makecell{發射星雲 \\ 反射星雲 \\ 反射星雲 \\ 暗星雲} & （四個一起） & \(\ang{;33;}\) & \(18^\text{h }18^\text{m }48^\text{s}\) & \(-\ang{19;45;3}\) \\ % RGB & > 6 hr \\
\end{longtable}
\end{center}

\subsection{武仙座 Hercules}
\begin{enumerate}
    \item 最亮星：武仙座\(\alpha\)星
    \item 全天第五大星座
\end{enumerate}
武仙座中的深空天體：
\newline

\begin{center}
\begin{longtable}{c | c | c | c | c | c}
    % \hline
    編號 & 種類 & 別名 & 視尺寸 & RA & Dec \\
    \hline
    \endfirsthead

    % \hline
    % 編號 & 種類 & 別名 & 視尺寸 & 拍攝技巧 & 曝光時間 \\
    % \hline
    % \endhead

    % \hline
    % \multicolumn{6}{r}{接下頁} \\
    \endfoot

    % \hline
    \endlastfoot

    M13/NGC 6205 & 球狀星團 & 武仙座球狀星團 & \(\ang{;20;} \times \ang{;40;}\) & \(16^\text{h }41^\text{m }42^\text{s}\) & \(+\ang{36;27;40}\) \\ % RGB & - \\
    M92/NGC 6341 & 球狀星團 & - & \(\ang{;10;}\) & \(17^\text{h }17^\text{m }07^\text{s}\) & \(+\ang{43;8;11}\) \\ % RGB & - \\
\end{longtable}
\end{center}

\subsection{巨蛇座 Serpens}
\begin{enumerate}
    \item 最亮星：巨蛇座\(\alpha\)星（天市右垣七），視星等2.63
\end{enumerate}
巨蛇座中的深空天體：
\newline

\begin{center}
\begin{longtable}{c | c | c | c | c | c}
    % \hline
    編號 & 種類 & 別名 & 視尺寸 & RA & Dec \\
    \hline
    \endfirsthead

    % \hline
    % 編號 & 種類 & 別名 & 視尺寸 & 拍攝技巧 & 曝光時間 \\
    % \hline
    % \endhead

    % \hline
    % \multicolumn{6}{r}{接下頁} \\
    \endfoot

    % \hline
    \endlastfoot

    M5/NGC 5904 & 球狀星團 & - & \(\ang{;23;}\) & \(15^\text{h }18^\text{m }33^\text{s}\) & \(+\ang{2;4;57}\) \\ % RGB & - \\
    M16/NGC 6611 & 發射星雲 & 星之皇后星雲/鷹星雲 & \(\ang{1;10;}\times\ang{;50;}\) & \(18^\text{h }18^\text{m }47^\text{s}\) & \(-\ang{13;48;25}\) \\ % RGB/SHO & > 4 hr \\
\end{longtable}
\end{center}

\subsection{蛇夫座 Ophiuchus}
\begin{enumerate}
    \item 最亮星：蛇夫座\(\alpha\)星（侯），視星等2.08
\end{enumerate}
蛇夫座中的深空天體：
\newline

\begin{center}
\begin{longtable}{c | c | c | c | c | c}
    % \hline
    編號 & 種類 & 別名 & 視尺寸 & RA & Dec \\
    \hline
    \endfirsthead

    % \hline
    % 編號 & 種類 & 別名 & 視尺寸 & 拍攝技巧 & 曝光時間 \\
    % \hline
    % \endhead

    % \hline
    % \multicolumn{6}{r}{接下頁} \\
    \endfoot

    % \hline
    \endlastfoot

    M9/NGC 6333 & 球狀星團 & - & \(\ang{;12;}\) & \(17^\text{h }19^\text{m }12^\text{s}\) & \(-\ang{18;30;59}\) \\ % RGB & - \\
    M10/NGC 6254 & 球狀星團 & - & \(\ang{;19;}\) & \(16^\text{h }57^\text{m }09^\text{s}\) & \(-\ang{4;5;57}\) \\ % RGB & - \\
    M12/NGC 6218 & 球狀星團 & - & \(\ang{;16;}\) & \(16^\text{h }41^\text{m }15^\text{s}\) & \(-\ang{1;56;15}\) \\ % RGB & - \\
    M14/NGC 6402 & 球狀星團 & - & \(\ang{;14;}\) & \(17^\text{h }37^\text{m }36^\text{s}\) & \(-\ang{3;14;44}\) \\ % RGB & - \\
    M19/NGC 6273 & 球狀星團 & - & \(\ang{;17;}\) & \(17^\text{h }02^\text{m }38^\text{s}\) & \(-\ang{26;16;5}\) \\ % RGB & - \\
    M62/NGC 6266 & 球狀星團 & - & \(\ang{;15;}\) & \(17^\text{h }01^\text{m }13^\text{s}\) & \(-\ang{30;6;45}\) \\ % RGB & - \\
    M107/NGC 6171 & 球狀星團 & - & \(\ang{;13;}\) & \(16^\text{h }32^\text{m }31^\text{s}\) & \(-\ang{13;3;12}\) \\ % RGB & - \\
    XSS J16271-2423 & 暗星雲 & 蛇夫座\(\rho\)星雲複合體 & \(\ang{6;30}\times\ang{4;30}\) & \(16^\text{h }28^\text{m }06^\text{s}\) & \(-\ang{24;32;30}\) \\ % RGB & > 4 hr
    LDN 1773/LDN 42 & 暗星雲 & 菸斗星雲 & \(\ang{6;;}\) & \(17^\text{h }27^\text{m }00^\text{s}\) & \(-\ang{26;56;0}\) \\
\end{longtable}
\end{center}

\subsection{天琴座 Lyra}
\begin{enumerate}
    \item 最亮星：織女一，視星等0.03，全天第五亮的星
\end{enumerate}
天琴座中的深空天體：
\newline

\begin{center}
\begin{longtable}{c | c | c | c | c | c}
    % \hline
    編號 & 種類 & 別名 & 視尺寸 & RA & Dec \\
    \hline
    \endfirsthead

    % \hline
    % 編號 & 種類 & 別名 & 視尺寸 & 拍攝技巧 & 曝光時間 \\
    % \hline
    % \endhead

    % \hline
    % \multicolumn{6}{r}{接下頁} \\
    \endfoot

    % \hline
    \endlastfoot

    M56/NGC 6779 & 球狀星團 & - & \(\ang{;8;48}\) & \(19^\text{h }16^\text{m }36^\text{s}\) & \(+\ang{30;11;4}\) \\ % RGB & - \\
    M57/NGC 6720 & 行星狀星雲 & 環狀星雲 & \(\ang{;3;50}\) & \(18^\text{h }53^\text{m }35^\text{s}\) & \(+\ang{33;1;42}\) \\ % RGB & - \\
\end{longtable}
\end{center}

\subsection{天鷹座 Aquila}
\begin{enumerate}
    \item 最亮星：河鼓二（天狼星），視星等0.8
\end{enumerate}
天鷹座中的深空天體：

\begin{center}
\begin{longtable}{c | c | c | c | c | c}
    % \hline
    編號 & 種類 & 別名 & 視尺寸 & RA & Dec \\
    \hline
    \endfirsthead

    % \hline
    % 編號 & 種類 & 別名 & 視尺寸 & 拍攝技巧 & 曝光時間 \\
    % \hline
    % \endhead

    % \hline
    % \multicolumn{6}{r}{接下頁} \\
    \endfoot

    % \hline
    \endlastfoot

    Barnard 142, 143 & 暗星雲 & \makecell{Barnard's E\\Nebula} & \(\ang{;30;}\) & \(19^\text{h }40^\text{m }42^\text{s}\) & \(+\ang{10;57;0}\) \\ % RGB & - \\
\end{longtable}
\end{center}
% \newpage

\subsection{天鵝座 Cygnus}
\begin{enumerate}
    \item 最亮星：天鵝座\(\alpha\)星（天津四），視星等1.3等
    \item 輦道增七：它位於天鵝的頭部，為一顆特別的雙星，兩顆星的顏色為一紅一藍
    \item 中心部分別名北十字星
\end{enumerate}
天鵝座中的深空天體：
\newline

\begin{center}
\begin{longtable}{c | c | c | c | c | c}
    % \hline
    編號 & 種類 & 別名 & 視尺寸 & RA & Dec \\
    \hline
    \endfirsthead

    % \hline
    % 編號 & 種類 & 別名 & 視尺寸 & 拍攝技巧 & 曝光時間 \\
    % \hline
    % \endhead

    % \hline
    % \multicolumn{6}{r}{接下頁} \\
    \endfoot

    % \hline
    \endlastfoot

    M29/NGC 6913 & 疏散星團 & - & \(\ang{;7;}\) & \(20^\text{h }23^\text{m }57^\text{s}\) & \(+\ang{38;30;27}\) \\ % RGB & - \\
    M39/NGC 7092 & 疏散星團 & - & \(\ang{;35;}\) & \(21^\text{h }31^\text{m }48^\text{s}\) & \(+\ang{48;26;17}\) \\ % RGB & - \\
    NGC 6826 & 行星狀星雲 & - & \(\ang{;;27}\times\ang{;;24}\) & \(19^\text{h }44^\text{m }48^\text{s}\) & \(+\ang{50;31;30}\) \\ % RGB & - \\
    C33/NGC 6992 & 發射星雲 & 東面紗星雲 & \(\ang{1;;}\times\ang{;30;}\) & \(19^\text{h }44^\text{m }48^\text{s}\) & \(+\ang{50;31;30}\) \\ % RGB/HOO & > 6 hr \\
    C34/NGC 6960 & 發射星雲 & 西面紗星雲 & \(\ang{1;10;}\times\ang{;6;}\) & \(20^\text{h }45^\text{m }57^\text{s}\) & \(+\ang{30;35;42}\) \\ % RGB/HOO & > 6 hr \\
    - & 發射星雲 & 東與西面紗星雲 & \(\ang{3;;}\) & - & - \\ % RGB/HOO & > 6 hr \\
    C20/NGC 7000 & 發射星雲 & 北美洲星雲 & \(\ang{2;;}\times\ang{1;40;}\) & \(20^\text{h }59^\text{m }17^\text{s}\) & \(+\ang{44;31;43}\) \\ % RGB/SHO & > 4 hr \\
    C27/NGC 6888 & WR星 & 眉月星雲 & \(\ang{;20;}\times\ang{;10;}\) & \(20^\text{h }12^\text{m }06^\text{s}\) & \(+\ang{38;21;17}\) \\ % RGB/HOO & > 10 hr \\
    IC 5070 & 發射星雲 & 鵜鶘星雲 & \(\ang{1;20;}\times\ang{1;10;}\) & \(20^\text{h }51^\text{m }01^\text{s}\) & \(+\ang{44;24;5}\) \\ % RGB/SHO & > 10 hr \\
    LDN 988 & 暗星雲 & - & \(\ang{3;;}\) & \(21^\text{h }08^\text{m }40^\text{s}\) & \(+\ang{48;57;17}\) \\ % RGB & > 15 hr \\
    Sh2-101 & 發射星雲 & 鬱金香星雲 & \(\ang{;16;}\times\ang{;9;}\) & \(19^\text{h }59^\text{m }53^\text{s}\) & \(+\ang{35;16;36}\) \\ % RGB/SHO & > 10 hr \\
    Barnard 147 & 暗星雲 & - & \(\ang{;11;}\) & \(20^\text{h }06^\text{m }59^\text{s}\) & \(+\ang{35;12;34}\) \\
\end{longtable}
\end{center}

% \newpage

\subsection{狐狸座 Vulpecula}
\begin{enumerate}
    \item 最亮星：狐狸座\(\alpha\)星（齊增五），視星等4.4等
\end{enumerate}
狐狸座中的深空天體：
\newline

\begin{center}
\begin{longtable}{c | c | c | c | c | c}
    % \hline
    編號 & 種類 & 別名 & 視尺寸 & RA & Dec \\
    \hline
    \endfirsthead

    % \hline
    % 編號 & 種類 & 別名 & 視尺寸 & 拍攝技巧 & 曝光時間 \\
    % \hline
    % \endhead

    % \hline
    % \multicolumn{6}{r}{接下頁} \\
    \endfoot

    % \hline
    \endlastfoot

    M27/NGC 6853 & 行星狀星雲 & 啞鈴星雲 & \(\ang{;8;}\times\ang{;5;48}\) & \(19^\text{h }59^\text{m }35^\text{s}\) & \(+\ang{22;43;15}\) \\ %RGB/HOO/SHO & > 8 hr \\
    Sh2-88 & 發射/反射/瀰漫星雲 & - & \(\ang{;18;}\) & \(19^\text{h }45^\text{m }59^\text{s}\) & \(+\ang{25;20;0}\) \\ %RGB & - \\
    Sh2-89 + Sh2-90 & 發射/反射星雲 & - & \(\ang{;4;42}\) & \(19^\text{h }50^\text{m }0^\text{s}\) & \(+\ang{26;40;0}\) \\ % RGB/SHO & - \\

\end{longtable}
\end{center}

\subsection{海豚座 Delphinus}
\begin{enumerate}
    \item 四顆主星排列成一個菱形
\end{enumerate}

\subsection{天秤座 Libra}
\begin{enumerate}
    \item 星座中最亮的四顆星\(\alpha\)、\(\beta\)、\(\gamma\)、\(\sigma\)構成一個四邊形
\end{enumerate}
天秤座中的深空天體：
\newline

\begin{center}
\begin{longtable}{c | c | c | c | c | c}
    % \hline
    編號 & 種類 & 別名 & 視尺寸 & RA & Dec \\
    \hline
    \endfirsthead

    % \hline
    % 編號 & 種類 & 別名 & 視尺寸 & 拍攝技巧 & 曝光時間 \\
    % \hline
    % \endhead

    % \hline
    % \multicolumn{6}{r}{接下頁} \\
    \endfoot

    % \hline
    \endlastfoot

    NGC 5897 & 球狀星團 & - & \(\ang{;12;48}\) & \(15^\text{h }17^\text{m }23^\text{s}\) & \(-\ang{21;00;36}\) \\ %RGB & - \\
    NGC 5898 & 橢圓星系 & - & \(\ang{;2;42}\times\ang{;2;36}\) & \(15^\text{h }18^\text{m }14^\text{s}\) & \(-\ang{24;05;53}\) \\ %RGB & - \\
    NGC 5728 & 棒旋星系 & - & \(\ang{;3;23}\) & \(14^\text{h }42^\text{m }24^\text{s}\) & \(-\ang{17;15;11}\) \\ %RGB & - \\
    NGC 5885 & 棒旋星系 & - & \(\ang{;3;30}\times\ang{;3;6}\) & \(15^\text{h }15^\text{m }03^\text{s}\) & \(-\ang{10;5;9}\) \\ %RGB & - \\
\end{longtable}
\end{center}

\subsection{盾牌座 Scutum}
\begin{enumerate}
    \item 最亮星：盾牌座\(\alpha\)，視星等3.85
\end{enumerate}
盾牌座中的深空天體：
\newline

\begin{center}
\begin{longtable}{c | c | c | c | c | c}
    % \hline
    編號 & 種類 & 別名 & 視尺寸 & RA & Dec \\
    \hline
    \endfirsthead

    % \hline
    % 編號 & 種類 & 別名 & 視尺寸 & 拍攝技巧 & 曝光時間 \\
    % \hline
    % \endhead

    % \hline
    % \multicolumn{6}{r}{接下頁} \\
    \endfoot

    % \hline
    \endlastfoot

    M11/NGC 6705 & 疏散星團 & 野鴨星團 & \(\ang{;22;48}\) & \(18^\text{h }51^\text{m }06^\text{s}\) & \(-\ang{6;16;11}\) \\ %RGB & - \\
    M26/NGC 6694 & 疏散星團 & - & \(\ang{;14;}\) & \(18^\text{h }45^\text{m }19^\text{s}\) & \(-\ang{9;23;0}\) \\ %RGB & - \\
\end{longtable}
\end{center}

\subsection{小熊座 Ursa Minor}
\begin{enumerate}
    \item 與大熊座的北斗七星相似的形狀，被稱為小北斗
    \item 斗柄處為小熊\(\alpha\)星（勾陳一），為現今的北極星。顏色為黃白色。是最亮的造父變星，其視星等變化範圍為1.97至2.00，是一個三合星系統，其超巨星主星擁有兩顆黃白主系星伴星
\end{enumerate}